\begin{textsample}{NEG Dim 4 – global\_south\_2023\_12 – Score 1.00 – global\_south\_2023\_12/104126956.txt}  \label{ex:f4_neg_320}
President Dr Arif Alvi, while questioning the world ' s conscience over the bloodshed in Gaza and unabated human rights abuses in occupied Kashmir, stated on Friday the world was committing the biggest mistake by not treating all humans equal.
The president, in his address during an event held in connection with International Human Rights Day here, said the way the innocent people, notably children, were being killed in Gaza would have a long-lasting impact on peace, as it would be difficult for the sufferers to forget such brutalities.
First Lady Samina Alvi, diplomats from various countries and representatives of international non-governmental organisations attended the event.
The president said that being a developing country, Pakistan was evolving with regards to human rights, but humanity still awaited the powerful nations to respect the human rights being violated in Palestine and Kashmir, where thousands of people had been killed, including women and children.
President Alvi called for a world order with no room for wars, adding that humanitarianism should override vested interests.
Instead of promoting the narrative of civilised @ @ @ @ @ @ @ @ @ @ no wars in the world, he went on to add.
The president said that Pakistan was a well-legislated country vis-a-vis human rights, backed by the constitution and human rights declaration.
It emphasised the will to ensure legal justice, financial equality, empowerment and mainstreaming of women as well as the differently-abled people, he added.
He told the gathering that in case of any human rights violation, Pakistan ' s state and judicial system took prompt action against those responsible for the rights violation and the people also discouraged such acts.
The president said the world had witnessed immense development, but deplored its inability to address the economic disparity.
He said that no religion allowed discrimination on the basis of class, creed and colour, regretting that this challenge was yet faced by the world.
Caretaker Federal Minister for Human Rights Khalil George said that Pakistan was playing an effective role in promotion of human rights and peace.
He said that Pakistani minorities carrying out their religious practices and celebrating their festivals manifested @ @ @ @ @ @ @ @ @ @ He told the gathering that the government had not only condemned the incidents like Jaranwala tragedy, but also brought those responsible to justice.
Contrarily, in the neighboring country of India, he added, Christian girls had been stripped naked and paraded around the city right under the nose of the state administration.
National of Commission for Human Rights Chairperson Rabia Javeria Agha said that despite the limited resources, the Commission processed 3,400 complaints and issued 12 reports on various issues related to women.

% matched lemmas: 
\end{textsample}
